\section{Related Work}
\label{sec:related-work}

\subsection{SPC reconstruction and accuracy assessment}
\label{subsec:spc-reconstruction-accuracy}

% 2.1改成SPC方法论述
% 2.2插一段事前拍摄的图片冗余分析
% 资源消耗的论述应该放到SPC后面，因为评价指标->发现资源浪费->观测规划方法研究
% 插一句承上启下的关系
% 2.2改成观测规划方法研究，不用强调通用性

SPC combines stereo constraints with brightness information to estimate surface topography and albedo. \citet{gaskell2008} established its use for small-body characterization and navigation through local terrain maps and global shape products. The mathematical formulation described by \citet{gaskell2023} couples the estimation of surface properties with refinements to observing parameters: stereo constrains the three-dimensional positions of surface features, while photoclinometry estimates local slopes and relative albedo from images acquired under different viewing and illumination conditions. The slopes are then integrated into elevations subject to geometric and topographic constraints. In practical mission workflows, an initial coarse shape supports image registration and the construction of local terrain patches, or maplets; repeated alignment, parameter adjustment, and terrain refinement incorporate new observations and support progressively finer models \citep{barnouin2020,palmer2022}. Reconstruction is therefore an iterative process whose outcome depends on both the available images and their use in successive model updates.Accuracy studies clarify how these updates can be evaluated. \citet{weirich2022} reconstructed a synthetic asteroid from simulated images and compared the recovered surface with known topography, examining the relationship between internal SPC diagnostics and absolute geometric error. For actual mission data, \citet{alasad2021} assessed successive Bennu models using maplet residuals, disagreement among overlapping maplets, image-based comparisons, and independent laser-altimeter measurements. These complementary assessments reveal regional deficiencies and track model development, while distinguishing internal consistency from agreement with the actual surface. In particular, agreement among overlapping maplets can remain high when they share a systematic error. Thus, model resolution and internal consistency alone do not establish geometric accuracy.Controlled sensitivity experiments further connect reconstruction quality to the input observations and processing procedure. \citet{palmer2024accuracy} varied imaging and processing conditions using synthetic terrain, quantified the resulting errors, and identified difficulties near steep slopes. The companion study by \citet{palmer2024geometry} examined observation combinations and reconstruction iterations, demonstrating the value of complementary topographic and albedo images and the diminishing benefit of additional processing iterations. Together, these studies provide a basis for interpreting reconstruction diagnostics and identifying conditions under which additional observations may improve a model.Recent work extends the estimation methods available for this process. \citet{driver2025} introduced Photoclinometry-from-Motion (PhoMo), combining keypoint-based structure from motion with a photometric factor graph that jointly estimates camera poses, landmark positions, Sun direction, surface normals, and albedos. Tests on Dawn imagery of Vesta and Ceres assess rendering quality and agreement with independently reconstructed terrain. AstroSplat incorporates planetary reflectance models into Gaussian splatting and evaluates image rendering alongside surface-normal agreement with PhoMo outputs \citep{nolan2026astrosplat}. These developments improve the automation and physical modeling of reconstruction from supplied images. Their evaluation also illustrates why photometric agreement and geometric accuracy require distinct evidence: a comparison with another reconstructed surface is subject to that reference's own errors. Connecting such reconstruction assessments to the selection of subsequent observations remains a separate planning task.

This paper connects this reconstruction and validation foundation to observation selection. After each observation round, SPC updates the surface model and regional accuracy indicators identify terrain that remains insufficiently constrained. These indicators then guide the next plan, whose images enter a further reconstruction and assessment round. The contribution is to use model evaluation as a continuing input to observation decisions. Known reference terrain in simulation provides the basis for checking whether improvements in the selected indicators correspond to reductions in surface error.

\subsection{Generalizable SPC observation planning and reconstruction feedback}
\label{subsec:spc-observation-planning}

General guidance for SPC observation planning follows from the complementary roles of stereo and photoclinometry. \citet{palmer2022} describe viewing and illumination conditions that support topography and albedo estimation, and \citet{palmer2024geometry} test combinations of observations that constrain these quantities. Such recommendations are applicable across targets, but translating them into observation decisions requires accounting for shape-dependent visibility, shadowing, and the observations already available for each region.\citet{piccinin2022} developed a generalizable imaging policy using deep reinforcement learning. Given prescribed spacecraft motion and pointing, the policy decides when to acquire images from the current observing geometry, observation history, mapping state, and remaining storage allowance. Its evaluation includes Eros, Itokawa, Bennu, and comet 67P, as well as changes in observing conditions and navigation uncertainty. The method demonstrates that a common policy can adapt observation timing across different targets while managing an image budget. Its mapping objective, however, combines observation-angle scores and image counts; the reported mapping gains do not directly quantify the surface error obtained by reconstructing SPC terrain from the selected images.\citet{wang2025} formulated observation planning from a low-resolution shape model as a combinatorial optimization problem. Their method combines a genetic algorithm with local search to select observing positions and imaging actions, accounting for illumination, visibility, repeated coverage, and observation costs. Noise added supports robust planning under position and pointing uncertainty. This provides a direct precedent for selecting observation combinations that satisfy SPC requirements with fewer positions and observations. As in the policy of \citet{piccinin2022}, reconstruction suitability is evaluated through prescribed geometric criteria and coverage statistics. Neither framework evaluates a newly reconstructed SPC model after an observation round and uses its regional quality diagnostics to determine the next round of observations.Active reconstruction research supplies complementary tools for representing the changing surface and selecting subsequent views. PC-NBV predicts candidate-view gains from a partial point cloud \citep{zeng2020}, SCONE estimates surface-coverage gains through occupancy and visibility prediction \citep{guedon2022}, and GenNBV learns observation policies that transfer across reconstruction environments \citep{chen2024}. PointNet's processing of unordered point sets \citep{qi2017pointnet} also motivates our modified PointNet representation of the reconstructed surface, from which continuous camera extrinsics are predicted within a reinforcement-learning planning framework.More directly, \citet{frahm2025vinnbv} predicts the relative reconstruction improvement of candidate viewpoints and uses that prediction for sequential view selection. Its training targets are obtained by reconstructing scenes with additional views and measuring the reduction in Chamfer distance to reference geometry. The policy supports limited observation budgets and transfers to unseen object categories. This is a close precedent for predicting observation utility from reconstruction outcomes, showing that quality-driven view selection already extends beyond coverage maximization. Its RGB-D reconstruction setting, however, does not address the coupled illumination, albedo, and terrain estimation required by SPC.

Our method extends this connection between reconstruction outcomes and observation decisions to iterative SPC reconstruction. Comparing regional quality indicators before and after each SPC update quantifies the benefit delivered by the acquired image set. With reinforcement learning based planning method, the policy associate observation choices with measured reconstruction improvement and prioritize complementary views of unresolved terrain under the available solar illumination as the reward. In addition to coverage or geometric mapping scores, The specific contribution is the coupling of this decision process to SPC reconstruction and regional accuracy assessment, with expected observation benefits checked against the improvements obtained after reconstruction.

\subsection{Mission image volumes, redundancy, and observation efficiency}
\label{subsec:mission-image-efficiency}

Small-body missions illustrate the scale of the image collections supporting surface reconstruction. \citet{gaskell2008} report that approximately 600 AMICA science images were analyzed during Hayabusa's characterization of Itokawa. For Rosetta, the SPC SHAP5 model of comet 67P used more than 5500 OSIRIS narrow-angle camera images acquired through June 2015 \citep{preusker2017}. In their operational account of OSIRIS-REx, \citet{palmer2022} report 2351, 5114, and 15,325 images for Bennu models v10, v20, and v42, respectively. These numbers describe image collections associated with particular reconstruction products, rather than total mission observation counts, and should not be summed across model versions. They nevertheless show that detailed modeling can involve thousands of images, with corresponding demands on storage, downlink, and processing.Recent mission design also illustrates tighter image budgets. \citet{ferrari2026navcam} allocate approximately 200 MB to compressed Milani NavCam asteroid images, corresponding to about 540 images planned for downlink over the mission. This is a prospective data budget serving navigation and science, rather than a completed SPC image collection. It highlights the need to select observations whose combined scientific value justifies their use of a limited downlink allocation.Repeated imaging serves several operational and scientific purposes. During Bennu Approach, uncertainty in target position and spacecraft pointing required $2\times2$, $3\times3$, and $4\times4$ mosaics to ensure coverage, while the Baseball Diamond campaign included seven passes to obtain observations for terrain reconstruction \citep{palmer2022}. Overlap provides margin against missed coverage and supplies complementary constraints. Its value nevertheless depends on what each image adds: once a region is sufficiently constrained, further observations with similar viewing and illumination conditions can consume resources while producing little improvement in its reconstructed geometry. Redundancy becomes avoidable waste for the modeling objective when additional observations provide negligible accuracy gains and displace observations of unresolved terrain. Large image counts alone do not establish how much of a mission's imaging was redundant, particularly when images also support navigation or other scientific objectives.Existing studies demonstrate the potential for more selective observation. Under controlled local-terrain test conditions, \citet{palmer2024geometry} obtained high-quality DTMs from five complementary images: four emphasizing topography and one emphasizing albedo. This result illustrates the importance of image complementarity, without prescribing a five-image budget for global mapping. At the planning level, one 67P simulation in \citet{piccinin2022} used 285 images with a learned policy instead of 500 with uniform scheduling, a 43\% reduction, while improving the geometric mapping score. \citet{wang2025} likewise reduce observing positions and imaging actions while maintaining prescribed geometric coverage. In robotic reconstruction, \citet{pan2025} combine a single next-best view with prediction of the remaining view set, improving coverage and reducing movement cost. These results motivate joint selection of complementary observations; quantifying observation savings for SPC requires comparing campaigns at matched reconstruction accuracy.\citet{zhang2025imagepairs} provide recent evidence from actual asteroid imagery. Their KD-tree method identifies overlapping image pairs in large, unordered Bennu, Vesta, and Ryugu data sets and supports selection under SPC observing constraints. In a local Vesta experiment, 70 candidate images were reduced to 22 overlapping images and then to five images with complementary viewing and illumination geometries. The selected images produced an SPC terrain model with 100 m grid spacing and a reported elevation RMSE of 61.16 m against a reference DEM. This connects selective image use to a reconstructed and evaluated product. However, the selection operates on an existing archive, after observation and downlink costs have already been incurred, and does not establish that the excluded images lacked value for other regions or objectives. Our planning problem moves this reasoning into successive observation rounds, using the current model to identify which new images are worth obtaining.

Our framework addresses this issue by using reconstruction feedback to concentrate new observations on regions with unresolved modeling needs. The marginal improvement measured after each round informs the value of further imaging, allowing the planner to reduce repeated observation and redirect resources toward useful complementary ones. Regional accuracy requirements also provide a stopping criterion once sufficient evidence supports the desired model quality. Fewer observations are useful when they satisfies the required reconstruction accuracy and respect the spacecraft's observing, storage, and communication limits. Evaluating image count and resource use together with achieved surface accuracy shows whether the proposed feedback reduces redundancy and wasted observing resources.
